% Adaptación del preámbulo de Paper III v1.5 español; no plantilla standalone de Pandoc.
\documentclass[english,11pt,a4paper]{article}
\usepackage[margin=1in]{geometry}
\usepackage{fontspec}
\setmainfont{lmroman10-regular.otf}[BoldFont=lmroman10-bold.otf,ItalicFont=lmroman10-italic.otf,BoldItalicFont=lmroman10-bolditalic.otf]
\setsansfont{lmsans10-regular.otf}[BoldFont=lmsans10-bold.otf]
\setmonofont{lmmono10-regular.otf}[Scale=MatchLowercase]
\newlength{\LeanCharWidth}
\usepackage{amsmath,amssymb}
\usepackage[english]{babel}
\usepackage{microtype,parskip}
\usepackage{graphicx,float,caption}
\usepackage{longtable,booktabs,array,calc}
\usepackage{xcolor}
\usepackage{xurl}
\usepackage{hyperref,bookmark}
\usepackage{etoolbox,needspace}
\setcounter{secnumdepth}{-\maxdimen}
\setlength{\emergencystretch}{3em}
\providecommand{\tightlist}{\setlength{\itemsep}{0pt}\setlength{\parskip}{0pt}}
\newcounter{none}
\AtBeginEnvironment{longtable}{\Needspace{8\baselineskip}\small}
\pretocmd{\section}{\Needspace{5\baselineskip}}{}{}
\pretocmd{\subsection}{\Needspace{4\baselineskip}}{}{}
\pretocmd{\subsubsection}{\Needspace{4\baselineskip}}{}{}
\captionsetup{font=small,labelfont=bf}
\urlstyle{same}
\hypersetup{colorlinks=true,linkcolor=blue,urlcolor=blue,citecolor=blue,pdflang={en-US}}
\title{Exact stability for nearly extremal chordal clique partitions}
\author{Juan Pablo Traverso Gianini\\
Independent researcher, Santiago, Chile}
\date{}\begin{document}\raggedbottom\maketitle
Series research note II, accompanying Papers IV and VI\\
Preprint 1.0 --- 18 October 2026

\subsection{Abstract}\label{abstract}

Let \(M(n)=\lfloor n(n+1)/6\rfloor\), and let
\(d=M(n)-\operatorname{cp}(G)\) be the deficit of a chordal graph. We
combine the structural theorem of Paper IV with the parametric local
formula of Paper VI. For sufficiently large order and
\(0\le d\le n/2000\), the vertices of degree greater than \(n/2\) form a
clique \(P\). Relative to this clique, the deficit is exactly the sum of
the loss caused by its size, the number of missing cross edges, and
twice the number of exterior edges. In particular, the graph differs
from the associated complete split graph in at most \(d\) edges, and an
optimal partition uses cliques of order at most four. The proof first
obtains a coarse edit bound from IV, uses it to enter the local regime
of VI, and then recovers the exact formula. This is an application of
the two papers, with a separate formal companion.

A clique partition covers every edge exactly once by cliques of order at
least two. Its minimum number of blocks is denoted by
\(\operatorname{cp}(G)\). We write \(\operatorname{cp}_{\le4}(G)\) when
blocks are restricted to order at most four. All graphs are finite and
simple. For disjoint vertex sets \(P,O\), the notation \(K_P\vee I_O\)
denotes the complete split graph with clique \(P\) and independent set
\(O\).

The distinction between a local formula and a structural theorem matters
here. Paper VI gives exact attainment once the graph is locally close to
a suitable complete split graph. Paper IV supplies that reference clique
for a nearly extremal chordal graph. The present note records how the
two statements combine and keeps the resulting deficit identity
explicit. The eventual Erdős 81 bound is already established in IV and
in the other developments compared in its Section 8. This note concerns
the stated stability consequence, not a new solution of that problem.

\subsection{1. Exact stability in a linear deficit
window}\label{exact-stability-in-a-linear-deficit-window}

Put

\[
M(n)=\left\lfloor\frac{n(n+1)}6\right\rfloor,
\qquad B(n,p)=p(n-p)-\binom p2,
\qquad \eta=10^{-16}.
\]

\Needspace{6\baselineskip}
\textbf{Theorem 1.1.} Using the uniform far-rounding and near-structure
theorems of Paper IV stated below, there exists an integer \(N\) such
that every chordal graph \(G\) of order \(n\ge N\), with deficit

\[
d=M(n)-\operatorname{cp}(G),\qquad 0\le d\le n/2000,
\]

has the following description. The set

\[
P=\{v:2d_G(v)>n\}
\]

is a clique. Write \(p=|P|\), \(O=V(G)\setminus P\), let \(A\) count
missing edges between \(P\) and \(O\), and put \(m=e(G[O])\). Then

\[
\boxed{\operatorname{cp}(G)=\operatorname{cp}_{\le4}(G)
=B(n,p)-A-2m,}
\tag{1}
\] \[
\boxed{d=\bigl(M(n)-B(n,p)\bigr)+A+2m,}
\tag{2}
\] \[
\boxed{|E(G)\triangle E(K_P\vee I_O)|=A+m\le d.}
\tag{3}
\]

In addition,

\[
(6p-2n-1)^2\le24d+25.
\tag{4}
\]

Equation (2) separates the deficit into quantities that can be read from
the graph once the clique is known. The first term measures the effect
of its size; the other terms count actual missing cross edges and
exterior edges. Equation (3) then bounds the total number of edits to
the associated complete split graph. The degree test identifies the same
clique without requiring a search over vertex bipartitions.

The size \(p\) need not be a maximizing size, such as
\(\lfloor(n+1)/3\rfloor\). This is why the parametric Proposition 2.1 of
VI, rather than its balanced Theorem A alone, is needed.

\subsection{2. Structural and rounding
inputs}\label{structural-and-rounding-inputs}

Let \(W^*(G)\) be the fractional maximum gain of a packing of triangles
and copies of \(K_4\), with gains 2 and 5. Write \(F_4(G)=e(G)-W^*(G)\).

\begin{enumerate}
\def\labelenumi{\arabic{enumi}.}
\tightlist
\item
  \textbf{Uniform rounding {[}1, Theorem 3.1{]}.} For \(\zeta=\eta/2\),
  once \(n\ge N(\eta/2)\), there is a physical packing whose completion
  into cliques has at most \[
  F_4(G)+(\eta/2)n^2
  \] blocks. This input applies uniformly to the graphs in question, not
  only to a preselected fractional optimizer.
\item
  \textbf{Near-structure witness {[}1, Proposition 4.3, Theorem 5.0 and
  Lemma 5.3{]}.} If \(n\ge4\cdot10^{12}\), \(G\) is chordal and \[
  F_4(G)\ge n^2/6-\eta n^2,
  \] IV produces a literal clique root \(P\), a physical packing and its
  completed partition of size \(L\), with literal exterior-edge and
  missing-spoke counts \(m,A\), satisfying \[
  L\le B(n,p)-m/20-A/2.
  \tag{5}
  \]
\end{enumerate}

In the first input, completing a packing means covering every unused
edge individually. The resulting object is a partition of the original
graph, so its number of blocks is an upper bound for
\(\operatorname{cp}(G)\). In the second input, the root, packing,
partition and counts belong to a single witness. We will use its
partition to control the counts, and its root to apply VI.

The counts in this witness are the actual exterior edges and missing
cross edges of its regularized clique root. Equation (5.15) of IV gives
(5); its stronger form (5.15a) concerns the same witness. This physical
identification is needed before the local theorem of VI can be applied.
The separate formal companion records the source interfaces.

\subsection{3. The order threshold}\label{the-order-threshold}

Choose \(N\) so that, for every \(n\ge N\),

\[
n\ge N(\eta/2),\qquad n\ge4\cdot10^{12},\qquad n\ge100,
\qquad \frac\eta2n^2>\frac n{2000}+1.
\tag{6}
\]

Such an \(N\) exists because \(\eta>0\). For the purely numerical last
condition, \(n\ge2\cdot10^{13}\) is already sufficient; the unknown
uniform-rounding threshold may be larger. This is an eventual theorem,
not an all-orders numerical assertion.

\subsection{4. Proof}\label{proof}

We first exclude the far regime. We then use the near witness to bound
the number of edited edges and the size of its clique root. Those two
bounds verify the hypotheses of VI; only after applying its exact
formula do we improve the edit bound and identify the root by degrees.

Suppose first that \(F_4(G)<n^2/6-\eta n^2\). Uniform rounding gives a
physical clique partition with fewer than

\[
\frac{n^2}{6}-\frac\eta2n^2
<\frac{n^2}{6}-\frac n{2000}-1
\le M(n)-d=\operatorname{cp}(G)
\]

blocks, a contradiction. The last weak inequality uses
\(M(n)\ge n^2/6-1\) and \(d\le n/2000\). Thus the near-structure input
applies.

Use its root \(P\), its actual counts \(m,A\) and its partition of size
\(L\). The square identity gives \(B(n,p)\le Q(n)=(2n+1)^2/24\). For
integer orders, \(Q(n)-M(n)\) is \(1/24\) when \(n\equiv0,2\pmod3\), and
\(9/24\) when \(n\equiv1\pmod3\). Thus \(Q(n)<M(n)+1\); since \(B(n,p)\)
is an integer, \(B(n,p)\le M(n)\). Hence

\[
M(n)-d=\operatorname{cp}(G)\le L
\le B(n,p)-m/20-A/2\le M(n)-m/20-A/2.
\]

The first two conclusions below follow by comparing the left and right
sides of this chain. For the third, multiply the bound \(m/20+A/2\le d\)
by 20 and use \(A\ge0\):

\[
B(n,p)\ge M(n)-d,\qquad m/20+A/2\le d,
\qquad m+A\le20d.
\tag{7}
\]

Because this root is a clique, the complete split graph
\(S=K_P\vee I_O\) differs from \(G\) in exactly \(m+A\) edges. Its
maximum modification degree therefore satisfies

\[
h=\Delta(G\triangle S)\le m+A\le20d\le n/100.
\tag{8}
\]

To control the root size, set \(Q(n)=(2n+1)^2/24\). Completing the
square gives

\[
Q(n)-B(n,p)=\frac{(6p-2n-1)^2}{24},
\qquad 0<Q(n)-M(n)<25/24.
\]

The square identity measures the distance of the root size from the
center of the split parabola. By (7), \(M(n)-B(n,p)\le d\); adding the
floor error \(Q(n)-M(n)\) and multiplying by 24 proves (4). The bound on
that floor error follows from \(Q(n)=n(n+1)/6+1/24\) and the defining
property of the floor. Since \(d\le n/2000\) and \(n\ge100\),

\[
24d+25\le0.012n+25\le n^2/100.
\]

Taking square roots gives \(|6p-2n-1|\le n/10\). Thus
\(1.9n+1\le6p\le2.1n+1\). For \(n\ge100\), these inequalities imply the
convenient coarse bounds

\[
3n/10\le p\le2n/5.
\tag{9}
\]

Put \(q=n-p\). Equations (8)--(9) yield

\[
q-p=n-2p\ge n/5\ge18h,
\qquad 4p-q=5p-n\ge n/2\ge36h.
\]

These are precisely

\[
p+18h\le q\le4p-36h,
\]

the hypotheses of VI Proposition 2.1 with its parameter
\(t_{\rm par}=4\). Indeed, its inequalities are
\(p+(4t_{\rm par}+2)h\le q\) and
\(q+t_{\rm par}(t_{\rm par}+5)h\le t_{\rm par}p\). The required
positivity of the core and the condition \(q\ge p\) follow from (9) and
the first inequality. The root has no missing internal edges. The local
formula is therefore \(pq-\binom p2-A-2m\), which is (1). Subtracting it
from \(M(n)\) gives (2). All three summands in (2) are nonnegative, so
\(A+m\le A+2m\le d\), proving (3). Notice that the initial bound \(20d\)
was only an entry budget; exact attainment improves it to \(d\).

It remains to check that the root just used is the set specified in the
theorem. A vertex of the clique has degree \(n-1\) in the complete split
reference and loses at most \(h\) incident edges. An exterior vertex has
reference degree \(p\) and gains at most \(h\). Thus, for \(a\in P\) and
\(x\in O\), the modification-degree bounds give

\[
d_G(a)\ge n-1-h\ge0.99n-1>n/2,
\qquad d_G(x)\le p+h\le0.41n<n/2.
\]

Thus this same root is exactly the degree-defined set in the statement.
This proves the theorem.

\subsection{5. Relation to the series}\label{relation-to-the-series}

The existence of the clique root and the exclusion of the far regime
come from Paper IV. The exact formula in the resulting neighborhood
comes from Paper VI. The initial bound \(A+m\le20d\) pays for entry into
that neighborhood; exact attainment then improves it to \(A+m\le d\).
The abstract transfer inequality of Paper V does not supply the
structural root. Its application to the eventual Erdős 81 bound is
developed in the separate IV--V note.

When \(d=0\), equation (2) forces \(A=m=0\) and \(B(n,p)=M(n)\). Thus
the graph is complete split, with a maximizing clique size. That
classification is already established in IV {[}1, Corollary 6.2{]} and
Okechukwu {[}4, Corollary 1.2{]}. The present note records an exact
deficit identity and degree recognition throughout the stated linear
window.

Okechukwu {[}4, Theorem 1.3{]} proves qualitative stability for graphs
with sublinear rooted simplicial defect; Section 7 asks for a
quantitative relation between the partition deficit and edit distance.
Paper IV {[}1, Theorem 6.1{]} already provides a linear edit bound to an
associated complete split graph and distinguishes that graph from one
with an optimizing core size. The present application combines IV with
VI to obtain the coefficient-one bound (3) and the exact identity (2) in
the smaller window \(d\le n/2000\). Its edit distance is to the
degree-identified graph \(K_P\vee I_O\), whose core size need not
maximize \(B(n,p)\). It therefore does not assert an optimal
quantitative rate for the entire class considered in {[}4{]}. The
original chordal clique-partition problem is due to Erdős, Ordman and
Zalcstein {[}5{]}.

\subsection{Formalization and
reproducibility}\label{formalization-and-reproducibility}

The separate Lean 4.35 project proves Theorem 1.1 from the proved inputs
of IV and VI. The formal companion identifies its final source freeze,
the public endpoint and physical adapters, the canonical root build and
axiom audits, and the pinned reproduction commands. Its selected
dependency chain excludes IV's final global-bound theorem while
retaining IV's structural and rounding constructors. The source and
audit records distinguish this project's eight modules from the exact
upstream source cuts.
\href{https://github.com/jtraverso/erdos-81-chordal-clique-partitions/tree/main/series_notes/NOTE_II_CHORDAL_STABILITY/05_formalization}{Sources
and formal companion}.

\subsection{Acknowledgments}\label{acknowledgments}

The author is deeply grateful to his wife María Paz and his children
Lucas, Juan Cristóbal, Francisca, Raimundo, and Benjamín for their love,
patience, and support.

\subsection{Use of artificial intelligence and computational
tools}\label{use-of-artificial-intelligence-and-computational-tools}

In the source papers used here, Claude, by Anthropic, and ChatGPT/Codex,
by OpenAI, were used to explore and check arguments and to prepare the
manuscripts. Aristotle, by Harmonic, contributed to exploring
approaches, searching for counterexamples, developing proofs, and
formalizing and reviewing them in Lean; its role was not limited to
translating finished proofs.

For this note, ChatGPT/Codex assisted in assembling and checking the
argument, developing the Lean interfaces, and preparing the manuscript
and audit materials. Aristotle contributed a separate formalization
proposal.

Certo and Jacobian were the computational tools used in the underlying
research program to explore instances, check identities, and assess
candidates, as disclosed in Paper IV. Their use is acknowledged
separately from the logical dependencies of the formal proof. The finite
regression checks performed specifically for this note use the
exact-arithmetic scripts identified in the accompanying audit.
Computations were reviewed within their stated scope; a finite search
does not replace a universal proof.

\Needspace{3\baselineskip}The author remains responsible for the arguments, citations, code, and
presentation. No AI system is listed as an author.

\subsection{References}\small\label{references}

{[}1{]} J. P. Traverso Gianini, \emph{Clique partitions with rooted
simplicial defect: quantitative stability and sharp eventual bounds},
Paper IV, preprint v1.23, October 1, 2026.
\href{https://github.com/jtraverso/erdos-81-chordal-clique-partitions/blob/054593443bdece587dda9686a411bb18b0ca34c5/preprints/PAPER_IV/01_manuscript/PAPER_IV_preprint_v1.23_en.md}{Manuscript
and materials}. Deposit:
\href{https://doi.org/10.5281/zenodo.23089131}{doi:10.5281/zenodo.23089131}.

{[}2{]} J. P. Traverso Gianini, \emph{Optimal clique partitions under
local perturbations}, Paper VI, Preprint 1.0. The proof uses the frozen
parametric source interface identified in the companion, October 18,
2026.
\href{https://github.com/jtraverso/erdos-81-chordal-clique-partitions/tree/main/preprints/PAPER_VI}{Manuscript
and materials}.

{[}3{]} J. P. Traverso Gianini, \emph{From realizability certificates to
the Erdős 81 bound}, separate research note accompanying Papers IV and
V, Preprint 1.0, October 18, 2026.
\href{https://github.com/jtraverso/erdos-81-chordal-clique-partitions/tree/main/series_notes/NOTE_I_ERDOS81_TRANSFER}{Manuscript
and materials}.

{[}4{]} O. Okechukwu, \emph{Clique partitions and bounded simplicial
defect}, arXiv:2609.20871v1 (2026).
\url{https://arxiv.org/abs/2609.20871v1}.

{[}5{]} P. Erdős, E. T. Ordman, and Y. Zalcstein, ``Clique partitions of
chordal graphs,'' \emph{Combinatorics, Probability and Computing}
\textbf{2} (1993), 409--415.
\url{https://doi.org/10.1017/S0963548300000808}.

\end{document}
